\section{The game and the rules dialect studied}
\label{sec:rules}

Canadian Fish, also called Literature, is played by six players in two teams of
three, seated so that the teams alternate. The 54-card deck is partitioned into
nine \emph{half-suits} of six cards, each player is dealt nine, and the team
that takes at least five half-suits wins. On their turn a player asks one
opponent for one named card; the ask is legal only if the named card lies in a
live half-suit, the asker holds at least one other card of that half-suit and
not the named card, and the target is an opponent who still holds cards. A hit
transfers the card and retains the turn, a miss passes the turn to the target,
and every ask, answer, transfer and card count is public. A half-suit is claimed
by a \emph{declaration}: an exact allocation naming, for each of the six cards,
which member of the declarer's own team holds it, awarded to the declaring team
if every assignment is correct and to the opposing team otherwise.

That is the whole game, and it is the same game the two preceding reports
studied. The dialect is unchanged from \vfour{} and \vfive{}, so this section
states the selection and points at the treatment rather than repeating it:
\S\ref{app:dialect} gives the clause table and the clauses this study's
mechanisms depend on, and the v0.4 report's Appendix A gives the dialect at the
length required to re-implement it \cite{fishbot04}.

\subsection{The dialect studied}
\label{sec:rules-dialect}

Published descriptions of the game document mutually incompatible variants
\cite{pagat}, so the dialect below is one explicit selection among them and not
a canonical rule set. It matches the strategy corpus this line of work draws on
\cite{develin} and the rules of record supplied by the project owner, and every
contested choice is an engine switch, so its effect is measurable rather than
assumed. The consequential choices are: a 54-card deck in nine half-suits, nine
cards per player; declaration legal at any moment, in or out of turn, by a seat
holding cards or not; a misdeclaration awarding the half-suit to the opposing
team; lowest seat index first among simultaneous declarers; an eight-level
willingness ladder in the forced endgame that arises when one whole team is
cardless; and a 400-ask action cap whose residue is adjudicated neutrally by
physical majority. No prearranged signalling conventions are available to any
policy, and team communication is limited to that willingness bit and to the
successor chosen by a cardless turn-holder. Table~\ref{tab:dialect} in
\S\ref{app:dialect} lists each choice against the switch that changes it.
\S\ref{sec:results-dialects} reports the primary result under
\vsixDialectN{} dialect perturbations; unless a result says otherwise, every
number in this report was produced under the settings above, and the action cap
was reached in \vsixLimitRate{}\% of games.

Two of those choices are places where the engine is \emph{more restrictive} than
the rules of record allow, and both are coordination channels a policy could in
principle use. The v0.5 study measured both, and this report inherits its
conclusions rather than re-opening them.

\paragraph{Turn transfer.} When the turn reaches a cardless seat whose teammates
both still hold cards, the rules of record permit the team to negotiate openly
over who receives it, sharing nothing but willingness. The engine instead has
the cardless seat choose unilaterally from its own belief, and solicits no
willingness bit. The channel is empty: a genuine multi-candidate transfer
decision arises \numTransferEvents{} times per game, and replacing the
unilateral rule with a \emph{ground-truth} chooser on one team, paired on common
deals, is worth $\numTransferDelta \pm \numTransferCI$ sets per game over
\numTransferGames{} paired games, with the two arms diverging at all in
\numTransferDiverged{}\% of games
(\filepath{research/v05/results/P8-coordination.md}, \S1.4). A legal mechanism
cannot beat the clairvoyant one, so that interval bounds the whole channel, and
no version of this policy has ever been built to exploit it.

\paragraph{Declaration arbitration.} The rules do not say who acts first when
two seats offer a declaration at the same instant, so the order is a modelling
choice. The engine scans by lowest seat index, deliberately declining to rank
candidates by stated confidence, because a confidence comparison moves a
statistic of one seat's hand to five others. The cost of that refusal is
\numArbCost{}~pp of win rate pooled over \numArbGames{} games, range
\numArbRange{}~pp across five seeds --- and a clairvoyant arbitrator, which
picks a proposer whose named allocation is actually correct, is worth
\numArbOracle{}~pp on the same design
(\filepath{research/v05/results/P6-verify-arbitration-cost.md}, \S\S3--4). The
information-safe rule therefore forgoes a third of a percentage point, and no
rule of any kind, safe or unsafe, informed or omniscient, recovers appreciably
more.

\subsection{What is unchanged from the v0.4 and v0.5 studies}
\label{sec:rules-unchanged}

Nothing in the dialect, the engine's rule implementation or the deal generator
changed between \vfive{} and \vsix{}. A reader who follows the citation to the
v0.4 report's rules treatment will find it accurate on every clause except two,
which no report has yet retracted and which \S\ref{sec:corr-dialect} corrects
here:

\begin{enumerate}
  \item \textbf{An undeclared half-suit does not score nothing.} The v0.4 report
        justifies the unconditional second stage of its forcing rule on the
        premise that an unclaimed half-suit is worthless. Under the dialect
        actually evaluated --- by both studies --- residue at the action cap is
        awarded by physical majority, and on a wrong post-horizon declaration
        the declaring team holds a mean \numResidueHeld{} of the six cards. The
        stage converts a probable award into a certain concession. The v0.4
        paper states the dialect correctly in its rules section and contradicts
        it two sections later.
  \item \textbf{Restarting the forced-endgame ladder sharpens nothing.} The v0.4
        report argues that restarting the willingness ladder after each
        declaration lets early, safer declarations resolve cards and so sharpen
        the allocations available for later ones. The argument is sound and
        inoperative: on every occasion a team went cardless with live half-suits
        --- \numLadderOccasions{} of them at seed \numLadderSeed{} --- exactly
        \numLadderLiveSets{} half-suit was live, so there is no ordering to
        exploit.
\end{enumerate}

Both are restated correctly in \S\ref{app:dialect-corrections} for a reader who
reaches the appendix first, and both bear on this study: the first is a premise
the declaration mechanism of \S\ref{sec:mech-declare} must not inherit, and the
second describes a structure that a more patient declaration policy would begin
to produce.
